% sesmesh — Solvent Excluded Surface meshes for molecules
% Copyright (c) 2026 Ugo Varetto
% SPDX-License-Identifier: MIT OR Apache-2.0

\section{The SES as a level set}
\label{sec:background}

\paragraph{Definitions.}
Let atom $i = 1,\dots,N$ have centre $\vect c_i$ and radius $r_i > 0$, let
$r_p \ge 0$ be the probe radius and $R_i = r_i + r_p$ the expanded radii.
The open SAS solid is $A = \bigcup_i B^\circ(\vect c_i, R_i)$ and the SAS is
$\partial A$. The \emph{probe-accessible region}
\begin{equation}
  P = \mathbb R^3 \setminus A
    = \{\vect x : \norm{\vect x - \vect c_i} \ge R_i \ \forall i\}
\end{equation}
is the closed set of feasible probe centres. The \emph{excluded region}
$E = \{\vect x : d(\vect x, P) \ge r_p\}$ is the complement of the union of
all probe balls that do not overlap an atom~\cite{greer1978}, and the SES is
$\partial E$.

\paragraph{The SES field.}
\sesmesh{} samples
\begin{equation}
  F(\vect x) = r_p - d(\vect x, P),
  \label{eq:field}
\end{equation}
which is negative inside the molecule, positive in the solvent and zero on
the SES. Quan and Stamm~\cite{quan2016} showed that this level set is
exactly the SES, singularities included: self-intersecting tori and
intersecting concave patches are features of the level set of $F$ and need
no special treatment. The algorithm relies on six elementary properties
(proofs are short and given in the design document of the
repository):
\begin{enumerate}[label=(P\arabic*)]
  \item $F$ is 1-Lipschitz, as every distance function is.
    \label{p:lipschitz}
  \item $F(\vect x) = r_p$, its maximum, for $\vect x \in P$: everything
    outside the SAS is trivially ``outside''.
  \item Inside $E$, $|F(\vect x)| = \operatorname{dist}(\vect x, \mathrm{SES})$
    exactly; outside $E$, $0 < F(\vect x) \le \operatorname{dist}(\vect x,
    \mathrm{SES})$. \label{p:dist}
  \item $F(\vect x) \le s_{\mathrm{vdW}}(\vect x) :=
    \min_i(\norm{\vect x - \vect c_i} - r_i)$, because
    $d(\vect x, P) \ge R_i - \norm{\vect x - \vect c_i}$ for every $i$: a
    cheap ``inside'' test. \label{p:vdw}
  \item Where the closest point $\vect p(\vect x) \in P$ is unique and
    $\vect x \notin P$, $\nabla F(\vect x) = (\vect p - \vect x) /
    \norm{\vect p - \vect x}$; on the SES the outward unit normal is
    $\vect n = (\vect p - \vect x)/r_p$. \label{p:normal}
  \item For $\vect x \in E$, the point $\vect x^* = \vect p + r_p(\vect x -
    \vect p)/\norm{\vect x - \vect p}$ lies on the SES: an exact projection
    from inside. \label{p:project}
\end{enumerate}
The SAS and vdW surfaces use $F_{\mathrm{SAS}}(\vect x) =
\min_i(\norm{\vect x - \vect c_i} - R_i)$ and $F_{\mathrm{vdW}}(\vect x) =
\min_i(\norm{\vect x - \vect c_i} - r_i)$, which are also 1-Lipschitz and
exact on their zero sets; an SES with $r_p = 0$ is the vdW surface.

\paragraph{Closest point on the SAS.}
For $\vect x \in A$ the closest point of $P$ lies on $\partial A$, a union of
spherical faces, circular arcs (pairwise sphere intersections) and vertices
(triple intersections). It lies in the relative interior of a face or an
arc, or is a vertex; on a sphere the only interior local minimum of the
distance is the radial projection, on a circle the planar projection. Hence
$d(\vect x, P)$ is the minimum over three candidate classes, each subject to
a validity test---the candidate point must not lie strictly inside another
expanded sphere (\cref{tab:classes}). Class S produces the convex patches
of the SES, class C the saddles and class V the concave patches, so the
winning class of a vertex is its patch type.

\begin{table}[t]
  \centering
  \caption{Candidate closest points on the SAS for a query point $\vect x$.
    For circle $(i,j)$ with $D = \norm{\vect c_j - \vect c_i}$,
    $|R_i - R_j| < D < R_i + R_j$: $\vect u = (\vect c_j - \vect c_i)/D$,
    $t = (D^2 + R_i^2 - R_j^2)/(2D)$, centre $\vect m = \vect c_i + t\vect u$,
    radius $\rho = (R_i^2 - t^2)^{1/2}$; $\vect y = \vect x - \vect m$,
    $a = \vect y \cdot \vect u$, $\vect w = \vect y - a\vect u$. $N(i)$ is the
    set of spheres intersecting sphere $i$ and $K(i,j)$ the set of spheres
    that cut circle $(i,j)$.}
  \label{tab:classes}
  \small
  \begin{tabular}{@{}llll@{}}
    \toprule
    Class (SES patch) & Candidate $\vect q$ & Distance & Valid iff \\
    \midrule
    S, sphere $i$ (convex) &
      $\vect c_i + R_i \frac{\vect x - \vect c_i}{\norm{\vect x - \vect c_i}}$ &
      $\bigl|\norm{\vect x - \vect c_i} - R_i\bigr|$ &
      $\norm{\vect q - \vect c_k} \ge R_k\ \forall k \in N(i)$ \\[4pt]
    C, circle $(i,j)$ (saddle) &
      $\vect m + \rho\, \vect w/\norm{\vect w}$ &
      $\bigl(a^2 + (\norm{\vect w} - \rho)^2\bigr)^{1/2}$ &
      $\norm{\vect q - \vect c_k} \ge R_k\ \forall k \in K(i,j)$ \\[4pt]
    V, vertex $(i,j,k)$ (concave) &
      precomputed $\vect v$ & $\norm{\vect x - \vect v}$ &
      precomputed (exposed only) \\
    \bottomrule
  \end{tabular}
\end{table}

Degenerate queries have deterministic answers: at $\vect x = \vect c_i$
class S is undefined and skipped (the point is deep inside), and for
$\norm{\vect w} \approx 0$ all points of a circle are equidistant, so the
minimum over an exposed arc is attained at an arc endpoint (a class-V
vertex) or, for a fully exposed circle, at a fixed point. Validity uses
$\norm{\vect q - \vect c_k}^2 \ge R_k^2 - \tau_k$; a wrong decision for a
candidate within $\varepsilon$ of a blocking sphere changes the computed
distance by $O(\varepsilon)$, because the true closest point is then on the
nearby arc or vertex. Tolerances are therefore accuracy parameters, not
topology hazards; the only ill-posed inputs are exactly fitting probes,
which a fixed tolerance decides deterministically (\cref{sec:arrangement}).

\paragraph{Cavities.}
$P$ may have bounded components: probe-accessible pockets that the solvent
cannot reach from outside. Each yields closed SES components whose outward
normals point into the enclosed solvent, so they have negative signed
volume, while external components have positive signed volume. This gives
an exact and cheap classification on the mesh (\cref{sec:cavities}).
